\paragraph{Proof of Proposition~\ref{prop:beta}.} Let $F$ be the (continuous) c.d.f.\ of the score and $U_i=F(S_i)$, i.i.d.\ uniform for i.i.d.\ data. With $k=\lceil(n+1)(1-\alpha)\rceil\le n$, $\hat q=S_{(k)}$ and $\mathrm{Cov}(\mathcal D_c)=\PP(S_{n+1}\le S_{(k)}\mid\mathcal D_c)=F(S_{(k)})=U_{(k)}\sim\mathrm{Beta}(k,n+1-k)$, and $n+1-k=\lfloor(n+1)\alpha\rfloor=l$. The mean is $k/(n+1)$ and the variance $k\,l/((n+1)^2(n+2))\approx\alpha(1-\alpha)/(n+2)$. For a bijection $\psi$ with $s'(\psi(h),\psi(y))=s(h,y)$, $s'(\psi(h),\psi(Y_i))=S_i$, so thresholds coincide and $\{u:s'(\psi(h),u)\le\hat q\}=\psi(\{u:s(h,u)\le\hat q\})$; $Y\in C\iff\psi(Y)\in\psi(C)$. The Aitchison score is $\|\mathrm{ilr}(y_\delta)-\mathrm{ilr}(\hh_\delta)\|_2$ since ilr is an isometry of clr.\hfill$\square$

\paragraph{Proof of Proposition~\ref{prop:measure}.} Let $V\in\mathbb R^{(K-1)\times K}$ have orthonormal rows spanning $\mathbf 1^\perp$, so $\mathrm{ilr}(p)=V\log p$. In the free coordinates $\tilde p=p_{1:K-1}$, $p=E\tilde p+e_K$ with $E=[I_{K-1};-\mathbf 1^\top]$, hence $\partial\log p/\partial\tilde p=P^{-1}E$ with $P=\mathrm{diag}(p)$ and the Jacobian of $\tilde p\mapsto\mathrm{ilr}(p)$ is $J=VP^{-1}E$. Let $Q=[V;\mathbf 1^\top/\sqrt K]$ (orthogonal) and $M=[E\mid p]$; $\det M=p_K+\mathbf 1^\top\tilde p=1$. Since $P^{-1}p=\mathbf 1$ and $V\mathbf 1=0$,
$QP^{-1}M=\begin{pmatrix}J&0\\ *&\sqrt K\end{pmatrix}$, so $\pm1/\prod_ip_i=\sqrt K\det J$, i.e.\ $|\det J|=(\sqrt K\prod_ip_i)^{-1}$ and $dA=|\det J|\,d\lambda$. (Numerically verified for $K=2,3,4,6$: relative error $<7\times10^{-10}$, Appendix Table~\ref{tab:jac}.)

(a) Let $R$ contain $\{p\in\Delta:\|p-p^0\|<\epsilon'\}$ with $p^0_i=0$, $p^0_j>0$ ($j\ne i$). On this set $\prod_{j\ne i}p_j$ is bounded below by $c>0$ for small $\epsilon'$, and $A(R)\ge c'\int_0^{\epsilon''}dp_i/p_i=\infty$. For any score continuous in $y$ on the closed simplex (all scores except Aitchison/DirNLL with $\delta=0$), the sublevel set $\{s<\hat q\}$ is relatively open, so it suffices that it contains one point in the relative interior of a facet; for $L_2$ this holds iff $\hat q>\hh_i\sqrt{K/(K-1)}$ for some $i$ (distance from $\hh$ to the facet).

(b) Write $c=\mathrm{clr}(p)$, $\hat c=\mathrm{clr}(\hh)$, so $p=\mathrm{softmax}(c)$. If $\|c-\hat c\|_2\le\rho$ then $|c_j-\hat c_j|\le\rho$ for all $j$ and $p_i/\hh_i=e^{c_i-\hat c_i}\sum_je^{\hat c_j}/\sum_je^{c_j}\le e^{\rho}e^{\rho}$. Hence $\prod_ip_i\le e^{2K\rho}\prod_i\hh_i$ on the ball, and $\lambda(B)=\int_B\sqrt K\prod p_i\,dA\le\sqrt Ke^{2K\rho}\prod_i\hh_i\,\omega_{K-1}\rho^{K-1}$, where the Aitchison ball is a Euclidean ball of radius $\rho$ in ilr coordinates with volume $\omega_{K-1}\rho^{K-1}$. Finally $d\mathrm{Dir}_a=c_a\prod p_i^{a-1}d\lambda=c_a\sqrt K\prod p_i^{a}\,dA$.\hfill$\square$

(Remark: the bound in (b) is informative only for moderate $\rho$. In our experiments $\hat q$ for the Aitchison score is large (e.g.\ 2.5 in Fig.~\ref{fig:tern}), and regions are large in the Lebesgue sense in the interior; we did not isolate the boundary regime empirically.)

\paragraph{Proof of Proposition~\ref{prop:gini}.} For $y\sim\mathrm{Dir}(\phi\mu)$, $\mathrm{Var}(y_i)=\mu_i(1-\mu_i)/(\phi+1)$ and $\mathrm{Cov}(y)=(\mathrm{diag}\mu-\mu\mu^\top)/(\phi+1)$, so $\EE\|y-\mu\|^2=\mathrm{tr}\,\mathrm{Cov}(y)=\sum_i\mu_i(1-\mu_i)/(\phi+1)=\Gini(\mu)/(\phi+1)$. For the multinomial, $\mathrm{Cov}(y)=(\mathrm{diag}\mu-\mu\mu^\top)/N$. Write $y=g/\sum g$ with $g_i\sim\mathrm{Gamma}(\phi\mu_i)$; the multivariate CLT for $g$ ($\mathrm{Var}\,g_i=\phi\mu_i$) and the delta method for the closure map give $\sqrt\phi(y-\mu)\Rightarrow\mathcal N(0,\mathrm{diag}\mu-\mu\mu^\top)$. By the continuous mapping theorem $(\phi+1)\|y-\mu\|^2\approx\sum_j\lambda_jZ_j^2$ with $\sum_j\lambda_j=\mathrm{tr}(\mathrm{diag}\mu-\mu\mu^\top)=\Gini(\mu)$. For a fixed pooled threshold the conditional coverage is therefore $\PP(\sum\lambda_jZ_j^2\le(\phi+1)\hat q^2)$; its argument's mean scale is $\Gini(\mu)$, which is the sense in which it is first-order governed by $\Gini(\mu)/(\phi+1)$. The ratio $\|y-\mu\|^2/\Gini(\mu)$ has mean $1/(\phi+1)$ independent of $\mu$; it is not distributionally constant because the eigenvalue profile $\lambda(\mu)$ changes the shape.\hfill$\square$

\paragraph{Proof of Proposition~\ref{prop:naive}.} Let $C$ be the realised coverage, $C-e\approx\mathcal N(0,\sigma_n^2)$ and, given $C$, $\hat c-C\approx\mathcal N(0,\tau^2)$ with $\tau^2=\alpha(1-\alpha)/T$ (using $C(1-C)\approx\alpha(1-\alpha)$). The Wald interval has half-width $\approx1.96\tau$ and excludes $e$ iff $|\hat c-e|>1.96\tau$, where $\hat c-e\approx\mathcal N(0,\sigma_n^2+\tau^2)$. With $\sigma_n^2/\tau^2=T/(n+2)$ the probability is $2\{1-\Phi(1.96\tau/\sqrt{\sigma_n^2+\tau^2})\}=2\{1-\Phi(1.96/\sqrt{1+T/(n+2)})\}$. An interval built from $R$ i.i.d.\ realisations of $C$ estimates $e$ with a standard error that includes $\sigma_n$, so has the nominal level up to the usual $t$-interval approximation.\hfill$\square$
